\section{Greek Words That Control the Argument}

The following are glosses from the registered Robinson--Pierpont Greek
artifact, not a new translation of Scripture. Its inline apparatus reports
selected Byzantine and Nestle--Aland differences; it is not itself the
Nestle--Aland edition.

\subsection{\emph{S\=emeion}: the work must be read}

The crowd follows because it sees Jesus's \emph{s\=emeia}, signs, yet Jesus says
it seeks him not because it saw signs but because it ate and was filled
(6:2, 26). The contrast is deliberate: optical sight is not yet perception.
In the chapter's narrative logic, the act must lead to the identity and
mission of Jesus. Neither reduction to social sharing nor reduction to a
self-contained sacramental symbol exhausts that Christological function.

\subsection{\emph{Ergazesthe} and \emph{pisteu\=ete}: work that receives}

The imperative ``work'' at 6:27 prompts a plural question about the works of
God. Jesus answers with a singular work: believing in the one God sent
(6:28--29). This sequence prevents Eucharistic interpretation from bypassing
faith. It also prevents faith from becoming a human achievement independent of
grace, since the same discourse says that coming requires the Father's drawing
and gift (6:44, 65).

\subsection{\emph{Al\=eth\=es}/\emph{al\=eth\=os}: true food}

At 6:55 the registered Byzantine text has adverbs: Christ's flesh is food
``truly'' and his blood drink ``truly.'' Its apparatus reports the
Nestle--Aland adjectives, ``true food'' and ``true drink.'' The grammatical
difference does not yield one reading literal and the other symbolic. Both
forms intensify the claim's truth; a critical edition is required before
making a text-critical argument from the difference.

\subsection{\emph{Phagein} and \emph{tr\=ogein}: eating without a magic lexicon}

John first uses forms of \emph{esthi\=o}/\emph{phagein}, the ordinary verb ``to
eat.'' From 6:54 onward he repeatedly uses \emph{tr\=ogein}. The latter can be
vividly glossed ``chew'' or ``munch,'' but it also functions as an ordinary
eating verb; John uses it at 13:18 when citing the one who eats bread with
Jesus. The shift contributes insistence and concreteness, especially after the
dispute in 6:52, but the Catholic case does not depend on claiming that one
Greek verb can only mean ``gnaw.'' Context---flesh, blood, true food, abiding,
and resurrection---does more work than an etymological slogan.

\subsection{\emph{Sarx}, \emph{pneuma}, and the ascending Son}

\emph{Sarx}, flesh, names the gift in 6:51--56 and then appears in the maxim
that flesh profits nothing (6:63). Between those uses Jesus asks what the
disciples will make of the Son of Man ascending where he was before (6:62).
The sequence relocates the scandal: the flesh to be given belongs to the Son
who descended and ascends. \emph{Pneuma}, Spirit, gives life; it does not turn
the incarnate flesh into a dispensable metaphor. The opposition is between
reception according to divine life and judgment confined to ordinary fleshly
possibility.

\subsection{\emph{Menein}: communion as abiding}

The verb ``abide'' at 6:56 belongs to a major Johannine field. The same Gospel
uses it for the Spirit remaining upon Jesus (1:32--33), disciples remaining in
Jesus's word (8:31), and branches abiding in the vine (15:1--10). Eucharistic
eating is thus personal and ecclesial communion, not ingestion isolated from
discipleship. The word does not dissolve sacrament into sentiment; it names the
relation the sacramental gift serves.

